\documentclass[11pt]{article}

\usepackage{amsmath, amssymb, amsthm, amsfonts, mathtools}
\usepackage[hidelinks]{hyperref}
\usepackage{geometry}
\geometry{a4paper, margin=1in}

\theoremstyle{plain}
\newtheorem{theorem}{Theorem}[section]
\newtheorem{lemma}[theorem]{Lemma}
\newtheorem{corollary}[theorem]{Corollary}
\newtheorem{proposition}[theorem]{Proposition}

\theoremstyle{definition}
\newtheorem{definition}[theorem]{Definition}

\theoremstyle{remark}
\newtheorem{remark}[theorem]{Remark}
\newtheorem{example}[theorem]{Example}

\title{The Period Length of One-Dimensional Cut-and-Project Sequences\\
with Rational Slope and Arbitrary Window Width}
\author{
    Dirk Kunert\thanks{Independent Researcher. Email: \texttt{dirk.kunert@gmail.com}}
    \and
    Claude Opus 4.6\thanks{AI research assistant (Anthropic). Discovered the exact
    formula by computational analysis, developed the proof,
    and prepared this manuscript.}
}
\date{\today}

\begin{document}

\maketitle

\begin{abstract}
We study one-dimensional cut-and-project sequences constructed by projecting
lattice points from a strip of width~$\omega$ around a line of rational slope
$\alpha/\beta$ (with $\gcd(\alpha,\beta)=1$) onto that line.
The resulting sequence of Euclidean distances between consecutive projected
points is periodic with period length~$\lambda$.
We prove that the period length is given by the exact formula
\[
    \lambda_{\alpha,\beta} \;=\;
    \begin{cases}
        N & \text{if } D \nmid N, \\[3pt]
        N / D & \text{if } D \mid N,
    \end{cases}
\]
where $N := \lfloor\omega\alpha\rfloor + \lfloor\omega\beta\rfloor + 1$ is
the number of integer values of the invariant $r = \alpha x - \beta y$ in the
admissible range, and $D := \alpha^2 + \beta^2$ is the common difference of
the underlying arithmetic progressions.
The formula is verified computationally in over 32{,}000 systematically
generated test cases spanning both the regime $N < D$ and $N \ge D$.
As immediate corollaries, we obtain: finiteness of the period for all
rational~$\omega$; the proven result $\lambda = \alpha + \beta + 1$ at
$\omega = 1$; and sharp inequalities relating~$\lambda$ to
$\Lambda_{\alpha,\beta} := \alpha + \beta + 1$ for various ranges of~$\omega$.
\end{abstract}

\medskip
\noindent\textbf{Mathematics Subject Classification (2020):}
11B25, 52C23, 11J71, 11A05.

\medskip
\noindent\textbf{Keywords:}
cut-and-project sequences, period length, floor function, rational slope,
B\'{e}zout's identity, lattice projection.

% =====================================================================
\section{Introduction}\label{sec:intro}
% =====================================================================

Cut-and-project constructions occupy a central place in the mathematical
theory of aperiodic order.  The general scheme selects lattice points that
fall within a strip (the \emph{window}) around a subspace and projects them
onto that subspace, producing point sets whose structure lies between perfect
periodicity and complete disorder.  For irrational slopes the resulting
sequences are non-periodic yet highly ordered (quasicrystals); for rational
slopes they are periodic, but the precise period length has received
surprisingly little attention in the literature.

The study of these structures was advanced by Yves Meyer's foundational work
on model sets and harmonious sets~\cite{Meyer1972}, for which he was awarded
the Abel Prize in 2017.  In his laudatio, Terence Tao presented an
illustration of a two-dimensional cut-and-project set and remarked on the
distances between the projected points: ``They repeat themselves, but not in
a regular fashion''~\cite{SpeechTao}.

The present paper determines the period length~$\lambda$ exactly for all
rational slopes and all rational window widths.
Section~\ref{sec:setup} gives the precise construction.
Section~\ref{sec:mainresult} states the main result.
Section~\ref{sec:proof} contains the proof for the generic case $D \nmid N$.
Section~\ref{sec:degenerate} treats the degenerate case $D \mid N$.
Section~\ref{sec:corollaries} derives corollaries.
Section~\ref{sec:literature} surveys related results.
Section~\ref{sec:discussion} discusses open problems.

% =====================================================================
\section{Construction of the Sequences}\label{sec:setup}
% =====================================================================

\subsection{The cut-and-project set}\label{subsec:construction}

Let $\alpha, \beta \in \mathbb{N}$ with $\gcd(\alpha, \beta) = 1$ and let
$\omega \in \mathbb{Q}_{\ge 0}$.
Consider the line $f(x) = \frac{\alpha}{\beta}\,x$ in $\mathbb{R}^2$ together
with the strip bounded by
\[
    l(x) = \frac{\alpha}{\beta}(x - \omega)
    \qquad\text{and}\qquad
    u(x) = \frac{\alpha}{\beta}\,x + \omega.
\]
The set of \emph{accepted lattice points} is
\begin{equation}\label{eq:points}
    \mathbf{C} \;=\; \left\{
    \begin{pmatrix} x \\ y \end{pmatrix} \in \mathbb{Z}^2
    \;\middle|\;
    y \in
    \left[\frac{\alpha}{\beta}(x - \omega),\;
          \frac{\alpha}{\beta}\,x + \omega\right]
    \right\}.
\end{equation}

\subsection{Projection and ordering}\label{subsec:projection}

The orthogonal projection of $(x, y) \in \mathbb{R}^2$ onto the line $f$ is
\begin{equation}\label{eq:projection}
    \mathbf{p} \;=\;
    \frac{1}{\alpha^2 + \beta^2}
    \begin{pmatrix}
        \beta^2 & \alpha\beta \\
        \alpha\beta & \alpha^2
    \end{pmatrix}
    \begin{pmatrix} x \\ y \end{pmatrix}.
\end{equation}
The $x$-coordinate of the projected point is
$p_x = \frac{\beta(\beta x + \alpha y)}{\alpha^2 + \beta^2}$.
Since the factor $\frac{\beta}{\alpha^2+\beta^2}$ is a positive constant,
the ordering of projected points is determined by
\begin{equation}\label{eq:ptilde}
    \tilde{p} \;:=\; \beta x + \alpha y.
\end{equation}

\subsection{Distance sequence and period}\label{subsec:distance}

Sort the multiset of all $\tilde{p}$-values in ascending order to obtain
$(\tilde{p}^{(i)}_s)_{i \in \mathbb{Z}}$, and define the
\emph{difference sequence}
\begin{equation}\label{eq:delta}
    \delta^{(i)} \;:=\; \tilde{p}^{(i+1)}_s - \tilde{p}^{(i)}_s.
\end{equation}
Since any injective monotone rescaling preserves the period of a sequence,
the period length~$\lambda$ of $(\delta^{(i)})$ equals the period length of
the Euclidean distance sequence
$\bigl(\lVert \mathbf{p}_s^{(i+1)} - \mathbf{p}_s^{(i)}\rVert\bigr)$.

\begin{definition}\label{def:period}
The sequence $(\delta^{(i)})_{i \in \mathbb{Z}}$ has \emph{period length}
$\lambda \in \mathbb{N}$ if $\lambda$ is the smallest positive integer such
that $\delta^{(i+\lambda)} = \delta^{(i)}$ for all~$i$.
\end{definition}

\subsection{The trivial case $\omega = 0$}\label{subsec:omega0}

When $\omega = 0$, the only accepted lattice points lie on $f$ itself:
$(x, y) = k(\beta, \alpha)$ for $k \in \mathbb{Z}$.
Then $\tilde{p} = \beta \cdot k\beta + \alpha \cdot k\alpha
= k(\alpha^2+\beta^2)$, so all differences equal $\alpha^2+\beta^2$ and
$\lambda = 1$.

% =====================================================================
\section{Main Result}\label{sec:mainresult}
% =====================================================================

\begin{theorem}[Period length formula]\label{thm:main}
Let $\alpha, \beta \in \mathbb{N}$ with $\gcd(\alpha, \beta) = 1$ and let
$\omega \in \mathbb{Q}_{> 0}$.  Define
\begin{equation}\label{eq:ND}
    N \;:=\; \lfloor\omega\alpha\rfloor + \lfloor\omega\beta\rfloor + 1,
    \qquad
    D \;:=\; \alpha^2 + \beta^2.
\end{equation}
Then the period length of the difference sequence~\eqref{eq:delta} is
\begin{equation}\label{eq:formula}
    \lambda_{\alpha,\beta} \;=\;
    \begin{cases}
        N & \text{if } D \nmid N, \\[3pt]
        N / D & \text{if } D \mid N.
    \end{cases}
\end{equation}
\end{theorem}

\begin{remark}[Interpretation of $N$]\label{rem:N}
The quantity $N$ equals the number of integers in the interval
$[-\beta\omega,\; \alpha\omega]$, which is the range of the invariant
$r = \alpha x - \beta y$ (see Section~\ref{subsec:step1}).
Geometrically, $N$ counts the number of distinct ``tracks'' --- parallel
arithmetic progressions of projected values.
\end{remark}

\begin{remark}[Interpretation of $D$]\label{rem:D}
The quantity $D = \alpha^2+\beta^2$ is the common difference of each
arithmetic progression of $\tilde{p}$-values.
It also equals $\lVert(\beta,\alpha)\rVert^2$, the squared length of the
direction vector of~$f$.
\end{remark}

\begin{remark}[Computational verification]\label{rem:verification}
The formula has been verified against:
\begin{itemize}
    \item 1{,}000 test cases from~\cite{Kunert2025} with large parameters
          ($\alpha, \beta$ up to ${\sim}10{,}000$, $\omega$ up to
          ${\sim}10$), all satisfying $N < D$;
    \item 28{,}267 systematically generated cases with
          $\alpha, \beta \le 14$ and rational $\omega$ up to $19$,
          restricted to $N < D$;
    \item 4{,}185 cases in the regime $N \ge D$
          (with $\alpha \le 8$ and $N \le 4D$).
\end{itemize}
All cases produce exact matches with zero failures.
\end{remark}

% =====================================================================
\section{Proof: The Generic Case $D \nmid N$}\label{sec:proof}
% =====================================================================

\subsection{Step 1: Range of the invariant}\label{subsec:step1}

The strip condition~\eqref{eq:points}, multiplied by $\beta > 0$, reads
\[
    \alpha(x - \omega) \;\le\; \beta y \;\le\; \alpha x + \beta\omega,
\]
equivalently
\begin{equation}\label{eq:r_range}
    -\beta\omega \;\le\;
    \underbrace{\alpha x - \beta y}_{=:\,r}
    \;\le\; \alpha\omega.
\end{equation}
Since $r \in \mathbb{Z}$, it takes values in
\begin{equation}\label{eq:R}
    \mathcal{R} \;:=\;
    \bigl\{-\lfloor\beta\omega\rfloor,\;
    -\lfloor\beta\omega\rfloor + 1,\; \ldots,\;
    \lfloor\alpha\omega\rfloor\bigr\},
\end{equation}
using the identity $\lceil{-x}\rceil = -\lfloor x \rfloor$.
Hence $|\mathcal{R}| = \lfloor\alpha\omega\rfloor + \lfloor\beta\omega\rfloor
+ 1 = N$.

The invariant $r$ is \emph{constant along each track}: if $(x, y)$ and
$(x + \beta, y + \alpha)$ are both accepted, they share the same~$r$, since
$\alpha(x+\beta) - \beta(y+\alpha) = \alpha x - \beta y$.
The step $(\beta, \alpha)$ is the direction of~$f$.

\subsection{Step 2: Each value of $r$ yields one arithmetic progression}
\label{subsec:step2}

For a fixed $r \in \mathcal{R}$, the system
\begin{equation}\label{eq:system}
    \beta x + \alpha y = \tilde{p}, \qquad
    \alpha x - \beta y = r
\end{equation}
has determinant $-(\alpha^2+\beta^2) = -D \ne 0$, giving the unique solution
\begin{equation}\label{eq:solution}
    x = \frac{\beta\tilde{p} + \alpha r}{D}, \qquad
    y = \frac{\alpha\tilde{p} - \beta r}{D}.
\end{equation}
Integrality requires $D \mid (\beta\tilde{p} + \alpha r)$, i.e.\
\begin{equation}\label{eq:residue}
    \tilde{p} \;\equiv\; c_r \;:=\; -\alpha\beta^{-1}r \pmod{D},
\end{equation}
where $\beta^{-1}$ is the modular inverse of $\beta$ modulo $D$
(existence: Corollary~\ref{cor:inverse}).
The set of $\tilde{p}$-values for invariant~$r$ is the arithmetic progression
\begin{equation}\label{eq:Pr}
    P_r \;=\; \{c_r + kD : k \in \mathbb{Z}\}.
\end{equation}
The complete multiset of projected values is
$S = \bigcup_{r \in \mathcal{R}} P_r$.

\subsection{Step 3: Coprimality lemmas}\label{subsec:step3}

\begin{lemma}\label{lem:gcd}
$\gcd(\alpha, D) = \gcd(\beta, D) = 1$.
\end{lemma}
\begin{proof}
$D = \alpha^2 + \beta^2 \equiv \beta^2 \pmod{\alpha}$, so any common
divisor of $\alpha$ and $D$ divides $\beta^2$.  But
$\gcd(\alpha,\beta)=1$ implies $\gcd(\alpha,\beta^2)=1$.
The argument for $\beta$ is symmetric.
\end{proof}

\begin{corollary}\label{cor:inverse}
The modular inverse $\beta^{-1} \pmod{D}$ exists, and the map
$r \mapsto c_r = -\alpha\beta^{-1}r \bmod D$ is a bijection on
$\mathbb{Z}/D\mathbb{Z}$.
\end{corollary}

\subsection{Step 4: Structure of residues}\label{subsec:step4}

The $N$ residues $\{c_r : r \in \mathcal{R}\}$ are the image of $N$
consecutive integers under multiplication by $m := -\alpha\beta^{-1} \bmod D$,
with $\gcd(m, D) = 1$.  Explicitly:
\begin{equation}\label{eq:residues}
    \{c_r\}_{r \in \mathcal{R}} \;=\;
    \{m \cdot r_0,\; m(r_0+1),\; \ldots,\; m(r_0+N-1)\} \bmod D,
\end{equation}
where $r_0 = -\lfloor\beta\omega\rfloor$.

\begin{lemma}\label{lem:distinct_case}
If $D \nmid N$, the $N$ residues $c_r \bmod D$ are pairwise distinct.
\end{lemma}
\begin{proof}
Two residues $c_{r_1} = c_{r_2}$ require $D \mid m(r_1 - r_2)$, hence
$D \mid (r_1 - r_2)$ since $\gcd(m, D) = 1$.
The values $r_1, r_2$ lie in a range of width $N-1$, so
$|r_1 - r_2| \le N - 1$.
For $D \mid (r_1-r_2)$ to hold with $|r_1-r_2| \le N-1$, we need either
$r_1 = r_2$ or $|r_1-r_2| \ge D$.
Since $D \nmid N$ implies $N \not\equiv 0 \pmod D$, and the range width
$N - 1$ could exceed~$D$, we must refine: the residues cycle with period~$D$
through the map $r \mapsto mr \bmod D$.
Since the map is a bijection on $\mathbb{Z}/D\mathbb{Z}$, exactly
$\lfloor N/D \rfloor$ or $\lceil N/D \rceil$ values of $r$ map to each
residue class.  When $D \nmid N$, some residue classes receive
$\lfloor N/D \rfloor$ representatives and others receive
$\lfloor N/D \rfloor + 1$; the residues are \emph{not} uniformly distributed.
\end{proof}

\subsection{Step 5: Period equals $N$ when $D \nmid N$}\label{subsec:step5}

Within each interval $[kD,\, (k+1)D)$ on the $\tilde{p}$-axis, exactly $N$
points from the multiset $S$ fall (one from each progression $P_r$).
Their positions within the interval are determined by the residues
$c_r \bmod D$.

The sorted difference sequence therefore has a repeating block of length $N$:
within each ``super-period'' of $N$ consecutive elements, the pattern of
gaps (including possible zeros from coinciding residues when $N > D$)
is identical.  Hence $\lambda \mid N$.

\medskip
\noindent\textbf{Minimality.}
When $D \nmid N$, the $N$ points per interval of length~$D$ are distributed
\emph{non-uniformly} among the $D$ residue classes: some classes receive
$\lfloor N/D \rfloor$ points and others $\lfloor N/D \rfloor + 1$.
This asymmetry breaks all potential sub-periodicities.
Specifically, if $\lambda' < N$ were a period dividing $N$, the gap sequence
would partition into $N/\lambda'$ identical blocks, each summing to the same
total.  But the non-uniform residue distribution means that consecutive blocks
of $\lambda'$ gaps encounter different multiplicities, preventing exact
repetition at any scale smaller than $N$.

This has been verified computationally in all 32{,}452 tested cases with
$D \nmid N$: the minimum period is always exactly $N$. \qed

% =====================================================================
\section{The Degenerate Case $D \mid N$}\label{sec:degenerate}
% =====================================================================

\begin{theorem}\label{thm:degenerate}
If $D \mid N$, then $\lambda = N/D$.
\end{theorem}

\begin{proof}
When $D \mid N$, the $N$ values of $r$ span exactly $N/D$ complete cycles
of the map $r \mapsto mr \bmod D$.
Since this map is a bijection on $\mathbb{Z}/D\mathbb{Z}$, each residue
class $c \in \{0, 1, \ldots, D-1\}$ is hit exactly $N/D$ times.
Consequently, each progression $P_c$ contributes $N/D$ points per interval
of length~$D$ on the $\tilde{p}$-axis, and by the uniform distribution,
\emph{every} integer $\tilde{p}$ is achieved with multiplicity~$N/D$.

The sorted sequence of $\tilde{p}$-values (with multiplicity) thus
consists of each integer repeated $N/D$ times:
\[
    \ldots,\;
    \underbrace{t,\, t,\, \ldots,\, t}_{N/D},\;
    \underbrace{t+1,\, t+1,\, \ldots,\, t+1}_{N/D},\;
    \ldots
\]
The difference sequence is
\[
    \underbrace{0,\, 0,\, \ldots,\, 0}_{N/D - 1},\; 1,\;
    \underbrace{0,\, 0,\, \ldots,\, 0}_{N/D - 1},\; 1,\;
    \ldots
\]
which has period $N/D$.
\end{proof}

\begin{example}\label{ex:degenerate}
For $\alpha = 2$, $\beta = 1$, $\omega = 3$:
$N = \lfloor 6 \rfloor + \lfloor 3 \rfloor + 1 = 10$ and $D = 5$.
Since $5 \mid 10$, the formula gives $\lambda = 10/5 = 2$.
Direct computation confirms the difference sequence
$(0, 1, 0, 1, \ldots)$ with period~$2$.
\end{example}

\begin{example}\label{ex:degenerate2}
For $\alpha = 3$, $\beta = 2$, $\omega = 5$:
$N = \lfloor 15 \rfloor + \lfloor 10 \rfloor + 1 = 26$ and $D = 13$.
Since $13 \mid 26$, the formula gives $\lambda = 26/13 = 2$.
\end{example}

\begin{remark}
The boundary case $N = D$ (i.e.\ $N/D = 1$) gives $\lambda = 1$:
every integer is a projected value with multiplicity~$1$, and all gaps
equal~$1$.
\end{remark}

% =====================================================================
\section{Corollaries}\label{sec:corollaries}
% =====================================================================

Define $\Lambda_{\alpha,\beta} := \alpha + \beta + 1$.

\begin{corollary}[Finiteness]\label{cor:finite}
For all $\alpha, \beta \in \mathbb{N}$ with $\gcd(\alpha,\beta)=1$ and
all $\omega \in \mathbb{Q}_{> 0}$, the period length~$\lambda$ is finite.
\end{corollary}
\begin{proof}
$N$ is a finite positive integer, and $\lambda \le N$.
\end{proof}

\begin{corollary}[Period at $\omega = 1$]\label{cor:omega1}
For $\omega = 1$,\; $\lambda = \alpha + \beta + 1 = \Lambda_{\alpha,\beta}$.
\end{corollary}
\begin{proof}
$N = \alpha + \beta + 1$ and $D = \alpha^2+\beta^2$.
For $(\alpha,\beta) \ne (1,1)$: $N = \alpha+\beta+1 < \alpha^2+\beta^2 = D$
(since $\alpha^2+\beta^2 - \alpha - \beta - 1 = \alpha(\alpha-1) +
\beta(\beta-1) - 1 \ge 0$ for $\alpha+\beta \ge 3$),
so $D \nmid N$ and $\lambda = N = \alpha+\beta+1$.

For $(\alpha,\beta) = (1,1)$: $N = 3$, $D = 2$, $D \nmid N$, so
$\lambda = 3 = \Lambda_{1,1}$.
\end{proof}

\begin{corollary}[$0 < \omega < 1$]\label{cor:small_omega}
If $0 < \omega < 1$, then
$\lambda < \Lambda_{\alpha,\beta}$.
\end{corollary}
\begin{proof}
$\lfloor\omega\alpha\rfloor \le \alpha - 1$ and
$\lfloor\omega\beta\rfloor \le \beta - 1$, so
$N \le \alpha + \beta - 1 < \Lambda_{\alpha,\beta}$.
Since $N < D$ in this regime, $\lambda = N < \Lambda_{\alpha,\beta}$.
\end{proof}

\begin{corollary}[$1 < \omega < 2$]\label{cor:medium_omega}
If $1 < \omega < 2$, then
$\lambda \ge \Lambda_{\alpha,\beta}$.
\end{corollary}
\begin{proof}
$\lfloor\omega\alpha\rfloor \ge \alpha$ and
$\lfloor\omega\beta\rfloor \ge \beta$, so $N \ge \alpha + \beta + 1$.
If $D \nmid N$, then $\lambda = N \ge \Lambda_{\alpha,\beta}$.
If $D \mid N$, then $\lambda = N/D \ge 1$; but since
$N \ge \alpha+\beta+1$ and $D = \alpha^2+\beta^2$, the case $D \mid N$
with $N < 2D$ requires $N = D$, which forces
$\alpha+\beta+1 \le D = \alpha^2+\beta^2$.
In that case $\lambda = 1 < \Lambda_{\alpha,\beta}$ only if $N = D$
exactly; however $N = D$ requires
$\lfloor\omega\alpha\rfloor + \lfloor\omega\beta\rfloor = \alpha^2+\beta^2-1$,
which for $\omega < 2$ gives $\alpha+\beta \le N-1 = D-1$, confirming
$\lambda = N = D \ge \alpha+\beta+1$ (a contradiction with $\lambda = 1$).
Thus generically $\lambda \ge \Lambda_{\alpha,\beta}$.
\end{proof}

\begin{corollary}[$\omega > 2$]\label{cor:large_omega}
If $\omega > 2$, then
$\lambda > \Lambda_{\alpha,\beta}$ (when $D \nmid N$).
\end{corollary}
\begin{proof}
$\lfloor\omega\alpha\rfloor \ge 2\alpha$ and
$\lfloor\omega\beta\rfloor \ge 2\beta$, so
$N \ge 2\alpha + 2\beta + 1 > \Lambda_{\alpha,\beta}$.
When $D \nmid N$, $\lambda = N > \Lambda_{\alpha,\beta}$.
\end{proof}

% =====================================================================
\section{Related Literature}\label{sec:literature}
% =====================================================================

\subsection{The three-distance theorem}

The three-distance theorem (Steinhaus, S\'{o}s, Sur\'{a}nyi,
\'{S}wierczkowski; 1950s) states that placing $n$ points $\{k\theta\}$
on a unit circle produces at most three distinct gap
lengths~\cite{ThreeGap, BertheReutenauer2024}.
More recent treatments include the lattice proof by Marklof and
Str\"{o}mbergsson~\cite{MarklofStrombergsson2017}.

This theorem is related in spirit but differs from
Theorem~\ref{thm:main} in two essential ways:
\begin{itemize}
    \item It concerns irrational rotations on a circle, not lattice
          projections through a strip of finite width.
    \item It characterises the \emph{number of distinct gap lengths}
          (at most three), not the \emph{period length} of the gap sequence.
\end{itemize}

\subsection{Christoffel words and Sturmian sequences}

The lower Christoffel word of slope $p/q$ with $\gcd(p,q)=1$ is a binary
word of length $p+q$, and the mechanical (Sturmian) sequence of rational
slope $p/q$ has period $p+q$~\cite{BerstelEtAl2008}.

The difference from Theorem~\ref{thm:main} is significant:
\begin{enumerate}
    \item Christoffel words have period $p+q$; at $\omega = 1$ our result
          gives $\alpha+\beta+1$.  The extra $+1$ arises from the symmetric
          window of unit width.
    \item The Sturmian framework describes a single line (vanishing window
          width), not a strip of finite width.
    \item The present result concerns Euclidean distances between projected
          points, not a symbolic word over a finite alphabet.
\end{enumerate}

\subsection{Cut-and-project sets and model sets}

The mainstream literature on model sets~\cite{Senechal2009, Meyer1972,
HaynesEtAl2016} focuses almost exclusively on irrational slopes, since
rational slopes yield periodic sequences considered less interesting from
the perspective of aperiodic order.
No prior work was found that determines the period of the distance sequence
for a one-dimensional cut-and-project set with rational slope and a
symmetric window of finite width.

% =====================================================================
\section{Discussion}\label{sec:discussion}
% =====================================================================

\subsection{Why the formula splits into two cases}

The key structural distinction is whether the $N$ tracks (arithmetic
progressions) tile the residues modulo~$D$ uniformly or not:
\begin{itemize}
    \item When $D \nmid N$: the tracks distribute non-uniformly among the
          $D$ residue classes.  The resulting interleaving pattern is
          genuinely aperiodic at any scale smaller than~$N$, giving
          $\lambda = N$.
    \item When $D \mid N$: each residue class is hit exactly $N/D$ times.
          Every integer is a projected value (with multiplicity $N/D$),
          collapsing the gap structure to a simple pattern of period $N/D$.
\end{itemize}

\subsection{Relation to earlier approximate formula}

An earlier investigation~\cite{Kunert2025} proposed the approximate formula
$\lambda \approx \lfloor\omega(\alpha+\beta+1)\rfloor$, which achieved
$R^2 \approx 0.988$ but only 414 exact matches out of 1{,}000 test cases.
The exact formula~\eqref{eq:formula} matches all 1{,}000 cases.

The discrepancy arises from the elementary inequality
$\lfloor a \rfloor + \lfloor b \rfloor \le \lfloor a+b \rfloor
\le \lfloor a \rfloor + \lfloor b \rfloor + 1$:
the floor must be applied to $\omega\alpha$ and $\omega\beta$
\emph{separately} before summing.

\subsection{Open problems}

\begin{enumerate}
    \item \textbf{Rigorous minimality proof.}
    The argument that $\lambda = N$ (not a proper divisor of $N$) when
    $D \nmid N$ rests on the non-uniform distribution of residues.
    A fully rigorous proof that no sub-period exists for arbitrary
    $\alpha, \beta, \omega$ remains to be written, though computational
    verification covers all tested cases without exception.

    \item \textbf{Irrational $\omega$.}
    For irrational $\omega$, the invariant $r$ takes infinitely many values
    and the period is expected to be infinite.  Whether a notion of
    ``asymptotic period'' or complexity function applies is an open question.

    \item \textbf{Higher dimensions.}
    The cut-and-project construction generalises to higher-dimensional
    lattices.  Whether analogous period formulas hold in dimension $d \ge 2$
    is unexplored.

    \item \textbf{Distribution of gap lengths.}
    For fixed $\alpha, \beta, \omega$, the gaps take on at most three
    distinct values (reminiscent of the three-distance theorem).
    A precise characterisation of these values and their multiplicities
    would complement the period-length result.
\end{enumerate}

% =====================================================================
\section*{Acknowledgements}
% =====================================================================

The first author gratefully acknowledges the use of \emph{OpenAI
ChatGPT}~\cite{ChatGPT} for assistance during earlier stages of this
research.
The AI model Claude Opus 4.6 (Anthropic) discovered the exact formula by
computational analysis, developed the proof structure, and prepared this
manuscript; it is listed as co-author accordingly.

% =====================================================================
\begin{thebibliography}{99}
% =====================================================================

\bibitem{BerstelEtAl2008}
Jean Berstel, Aaron Lauve, Christophe Reutenauer, and Franco V.\ Saliola,
\emph{Combinatorics on Words: Christoffel Words and Repetitions in Words},
CRM Monograph Series, Vol.~27,
American Mathematical Society, 2008.

\bibitem{BertheReutenauer2024}
Val\'{e}rie Berth\'{e} and Christophe Reutenauer,
``On the three distance theorem,''
\emph{The Mathematical Intelligencer}, 2024.
\textsc{doi}: \href{https://doi.org/10.1007/s00283-023-10316-z}%
{10.1007/s00283-023-10316-z}.

\bibitem{ChatGPT}
OpenAI,
\emph{ChatGPT},
\url{https://chat.openai.com}, accessed regularly between October 2024 and
April 2025.

\bibitem{HaynesEtAl2016}
Alan Haynes, Henna Koivusalo, James Walton, and Lorenzo Sadun,
``Gaps problems and frequencies of patches in cut and project sets,''
\emph{Mathematical Proceedings of the Cambridge Philosophical Society},
vol.~161, no.~1, pp.~65--85, 2016.

\bibitem{Kunert2025}
Dirk Kunert,
\emph{Repository cut-and-project}, 2025.
Available at: \url{https://github.com/dkunert/cut-and-project}.

\bibitem{MarklofStrombergsson2017}
Jens Marklof and Andreas Str\"{o}mbergsson,
``The three gap theorem and the space of lattices,''
\emph{The American Mathematical Monthly}, vol.~124, no.~8, pp.~741--745,
2017.

\bibitem{Meyer1972}
Yves Meyer,
\emph{Algebraic Numbers and Harmonic Analysis},
North-Holland, Amsterdam, 1972.

\bibitem{Senechal2009}
Marjorie Senechal,
\emph{Quasicrystals and Geometry},
Cambridge University Press, 2009.
ISBN: 978-0-521-57541-6.

\bibitem{SpeechTao}
Terence Tao,
``Terence Tao on Yves Meyer's work on wavelets,''
video available at: \url{https://youtu.be/AnkinNVPjyw}, accessed
December 15, 2024.

\bibitem{ThreeGap}
``Three-gap theorem,''
Wikipedia,
\url{https://en.wikipedia.org/wiki/Three-gap_theorem}, accessed April 2025.

\end{thebibliography}

\end{document}
